% Generated draft: requires mathematical and source review.
\documentclass[11pt]{article}
\usepackage[margin=25mm]{geometry}
\usepackage{fontspec}
\setmainfont{Latin Modern Roman}
\usepackage{amsmath,amssymb,mathtools,unicode-math}
\setmathfont{Latin Modern Math}
\usepackage{enumitem,needspace,xurl,hyperref}
\setlength{\parindent}{0pt}
\setlength{\parskip}{5pt}
\setlength{\emergencystretch}{3em}
\newcommand{\N}{\mathbb{N}}
\newcommand{\Z}{\mathbb{Z}}
\newcommand{\Q}{\mathbb{Q}}
\newcommand{\R}{\mathbb{R}}
\newcommand{\C}{\mathbb{C}}
\newcommand{\F}{\mathbb{F}}
\newcommand{\A}{\mathbb{A}}
\newcommand{\PP}{\mathbb{P}}
\newcommand{\E}{\mathbb{E}}
\newcommand{\eps}{\varepsilon}
\newcommand{\dd}{\,\mathrm{d}}
\begin{document}
\section*{UnsolvedMath ID: 30001255}
\textbf{Source problem code:} catalogue entry
\subsection{Definitions and mathematical statement}
Let \(F\) be a field with \(\operatorname{char}(F)=2\). A symmetric bilinear form on an \(F\)-vector space \(V\) is a map \(b\colon V\times V\to F\) that is linear in each argument and satisfies \(b(v,w)=b(w,v)\). Its radical is
\[
\operatorname{rad}(b)=\{v\in V: b(v,w)=0\text{ for every }w\in V\}.
\]
In the usual finite-dimensional setting, regularity means \(\operatorname{rad}(b)=\{0\}\). The form is anisotropic if \(b(v,v)\ne 0\) for every \(v\ne 0\). The supplied record does not explicitly state the dimension convention.

For each field \(F\) of characteristic \(2\), each integer \(h>0\), and each integer \(d\geq 0\), the source asks: which anisotropic regular symmetric bilinear forms \(b\) over \(F\) satisfy
\[
\operatorname{height}(b)=h,
\qquad
\operatorname{degree}(b)=d?
\]
Here \(\operatorname{height}(b)\) and \(\operatorname{degree}(b)\) denote the invariants intended in the cited report. Their definitions and conventions are absent from the supplied record, so the displayed conditions cannot be made fully explicit from this evidence alone. The record also does not specify whether the requested classification is up to isometry, similarity, or another equivalence relation. These qualifications require resolution from the source before this becomes a complete mathematical formulation.
\subsection{Short English statement}
For every field of characteristic \(2\) and every pair of integers \(h>0\), \(d\geq 0\), classify anisotropic regular symmetric bilinear forms of height \(h\) and degree \(d\). The supplied record leaves the definitions of height and degree, the dimension convention, and the classification equivalence unspecified.
\subsection{Sources}
\begin{itemize}
\item Supplied local UnsolvedMath registry record, permanent ID \(30001255\), title: Anisotropic Bilinear Forms in Characteristic Two; exact source problem code: \texttt{catalogue entry}. The record gives the proposed year as 2009 and supplies no individual proposer.
\item \emph{Quadratic Forms and Linear Algebraic Groups} (2009), Oberwolfach Reports, pp.\ 1375--1426, DOI 10.4171/owr/2009/25. Supplied source URL and literature-source link, labelled ``original Oberwolfach report'': \url{https://doi.org/10.4171/owr/2009/25}.
\item Literature assessment checked at \texttt{2026-08-21}: supplied dated curation, not independently verified here. The supplied status is \texttt{partially\_solved}. The assessment says that the report completely resolves a displayed height-two exceptional case and several low-degree cases, while Question 1 asks for all heights and degrees and is not thereby classified in full. The supplied background explicitly requests independent verification of time-sensitive status and cited claims.
\item The registry labels the statement \texttt{corrected\_verified} and reports that its mathematical content was checked against the official report, with notation, context, and question form normalized without changing the source-level claim. This is a supplied verification claim, not verification performed for this draft.
\end{itemize}
\end{document}
